% ATIVIDADES DE SONS — Volume 2 (gerado por expandir_volume.py)

\chapter{Atividades de Cada Som}

No 2º ano, cada atividade trabalha um som-alvo com foco em consciência fonológica, articulação (Boquinha) e grafia. Cada atividade é uma plancheta reproduzível de página inteira; guarde as planchetas feitas: elas mostram o progresso.

%% ATIVIDADE 1: FONEMA CH EM POSIÇÃO INICIAL

\plancheta{Atividade de Som 1 --- Fonema CH em posição inicial}

\subsection{Atividade de Som 1 --- Fonema CH em posição inicial}
\begin{exercicio}[Atividade de Som 1 --- Fonema CH em posição inicial]

\metadata{Som01}{Silábico-Alfabético}{EF02LP01, EF02LP04}{D1}{EF02LP04}{CNE/CEB nº 2/2012}

\kumontiming{2}{90}

\noindent\textbf{Nome:} \underline{\hspace{3.2cm}} \quad \textbf{Data:} \underline{\hspace{1.6cm}} \quad \textbf{Nota:} \underline{\hspace{0.9cm}}/10


\noindent\textbf{Comando:} Peça que o aluno fale devagar: CHA, CHE, CHI, CHO, CHU. Circule todas as figuras que começam com o som de CH.


\begin{center}
\begin{tikzpicture}
  \draw[thick, dashed] (0,0) rectangle (12,6);
\end{tikzpicture}
\end{center}

\noindent\textbf{Escreva 2 palavras com o som da atividade:}

\begin{center}
  \underline{\hspace{3cm}} \quad \underline{\hspace{3cm}}
\end{center}

\noindent\textbf{Desenhe uma situação com a palavra da atividade:}

\begin{center}
\begin{tikzpicture}
  \draw[thick, dashed] (0,0) rectangle (12,4.5);
\end{tikzpicture}
\end{center}

\end{exercicio}

%% ATIVIDADE 2: FONEMA CH EM POSIÇÃO MEDIAL

\plancheta{Atividade de Som 2 --- Fonema CH em posição medial}

\subsection{Atividade de Som 2 --- Fonema CH em posição medial}
\begin{exercicio}[Atividade de Som 2 --- Fonema CH em posição medial]

\metadata{Som02}{Silábico-Alfabético}{EF02LP01, EF02LP04}{D1}{EF02LP04}{CNE/CEB nº 2/2012}

\kumontiming{2}{90}

\noindent\textbf{Nome:} \underline{\hspace{3.2cm}} \quad \textbf{Data:} \underline{\hspace{1.6cm}} \quad \textbf{Nota:} \underline{\hspace{0.9cm}}/10


\noindent\textbf{Comando:} Diga a palavra CHUVEIRO e peça que o aluno bata palmas a cada sílaba: CHU-VEI-RO. Agora encontre outras palavras com CH no meio.


\begin{center}
\begin{tikzpicture}
  \draw[thick, dashed] (0,0) rectangle (12,6);
\end{tikzpicture}
\end{center}

\noindent\textbf{Escreva 2 palavras com o som da atividade:}

\begin{center}
  \underline{\hspace{3cm}} \quad \underline{\hspace{3cm}}
\end{center}

\noindent\textbf{Desenhe uma situação com a palavra da atividade:}

\begin{center}
\begin{tikzpicture}
  \draw[thick, dashed] (0,0) rectangle (12,4.5);
\end{tikzpicture}
\end{center}

\end{exercicio}

%% ATIVIDADE 3: FONEMA LH EM POSIÇÃO MEDIAL

\plancheta{Atividade de Som 3 --- Fonema LH em posição medial}

\subsection{Atividade de Som 3 --- Fonema LH em posição medial}
\begin{exercicio}[Atividade de Som 3 --- Fonema LH em posição medial]

\metadata{Som03}{Silábico-Alfabético}{EF02LP01, EF02LP04}{D1}{EF02LP04}{CNE/CEB nº 2/2012}

\kumontiming{2}{90}

\noindent\textbf{Nome:} \underline{\hspace{3.2cm}} \quad \textbf{Data:} \underline{\hspace{1.6cm}} \quad \textbf{Nota:} \underline{\hspace{0.9cm}}/10


\noindent\textbf{Comando:} Diga COELHO e peça que o aluno sinta a língua no céu da boca. Escreva palavras com LH que tenham o som lateral.


\begin{center}
\begin{tikzpicture}
  \draw[thick, dashed] (0,0) rectangle (12,6);
\end{tikzpicture}
\end{center}

\noindent\textbf{Escreva 2 palavras com o som da atividade:}

\begin{center}
  \underline{\hspace{3cm}} \quad \underline{\hspace{3cm}}
\end{center}

\noindent\textbf{Desenhe uma situação com a palavra da atividade:}

\begin{center}
\begin{tikzpicture}
  \draw[thick, dashed] (0,0) rectangle (12,4.5);
\end{tikzpicture}
\end{center}

\end{exercicio}

%% ATIVIDADE 4: FONEMA NH EM POSIÇÃO MEDIAL

\plancheta{Atividade de Som 4 --- Fonema NH em posição medial}

\subsection{Atividade de Som 4 --- Fonema NH em posição medial}
\begin{exercicio}[Atividade de Som 4 --- Fonema NH em posição medial]

\metadata{Som04}{Silábico-Alfabético}{EF02LP01, EF02LP04}{D1}{EF02LP04}{CNE/CEB nº 2/2012}

\kumontiming{2}{90}

\noindent\textbf{Nome:} \underline{\hspace{3.2cm}} \quad \textbf{Data:} \underline{\hspace{1.6cm}} \quad \textbf{Nota:} \underline{\hspace{0.9cm}}/10


\noindent\textbf{Comando:} Diga NINHO e peça que o aluno feche a boca ao final (som nasal). Liste palavras com NH que tenham este gesto.


\begin{center}
\begin{tikzpicture}
  \draw[thick, dashed] (0,0) rectangle (12,6);
\end{tikzpicture}
\end{center}

\noindent\textbf{Escreva 2 palavras com o som da atividade:}

\begin{center}
  \underline{\hspace{3cm}} \quad \underline{\hspace{3cm}}
\end{center}

\noindent\textbf{Desenhe uma situação com a palavra da atividade:}

\begin{center}
\begin{tikzpicture}
  \draw[thick, dashed] (0,0) rectangle (12,4.5);
\end{tikzpicture}
\end{center}

\end{exercicio}

%% ATIVIDADE 5: CONTRASTE NH × LH

\plancheta{Atividade de Som 5 --- Contraste NH × LH}

\subsection{Atividade de Som 5 --- Contraste NH × LH}
\begin{exercicio}[Atividade de Som 5 --- Contraste NH × LH]

\metadata{Som05}{Silábico-Alfabético}{EF02LP01, EF02LP04}{D1}{EF02LP04}{CNE/CEB nº 2/2012}

\kumontiming{2}{90}

\noindent\textbf{Nome:} \underline{\hspace{3.2cm}} \quad \textbf{Data:} \underline{\hspace{1.6cm}} \quad \textbf{Nota:} \underline{\hspace{0.9cm}}/10


\noindent\textbf{Comando:} Compare NINHO (ar sai pelo nariz) com MILHO (ar sai lateral). Diga as duas palavras e peça que o aluno diga qual é qual.


\begin{center}
\begin{tikzpicture}
  \draw[thick, dashed] (0,0) rectangle (12,6);
\end{tikzpicture}
\end{center}

\noindent\textbf{Escreva 2 palavras com o som da atividade:}

\begin{center}
  \underline{\hspace{3cm}} \quad \underline{\hspace{3cm}}
\end{center}

\noindent\textbf{Desenhe uma situação com a palavra da atividade:}

\begin{center}
\begin{tikzpicture}
  \draw[thick, dashed] (0,0) rectangle (12,4.5);
\end{tikzpicture}
\end{center}

\end{exercicio}

%% ATIVIDADE 6: QUE E QUI: O U MUDO

\plancheta{Atividade de Som 6 --- QUE e QUI: o U mudo}

\subsection{Atividade de Som 6 --- QUE e QUI: o U mudo}
\begin{exercicio}[Atividade de Som 6 --- QUE e QUI: o U mudo]

\metadata{Som06}{Silábico-Alfabético}{EF02LP01, EF02LP04}{D1}{EF02LP04}{CNE/CEB nº 2/2012}

\kumontiming{2}{90}

\noindent\textbf{Nome:} \underline{\hspace{3.2cm}} \quad \textbf{Data:} \underline{\hspace{1.6cm}} \quad \textbf{Nota:} \underline{\hspace{0.9cm}}/10


\noindent\textbf{Comando:} Diga QUEIJO sem dizer o U: QUE. Compare com QUERO (QUE) e QUILO (QUI). A criança deve perceber que o U não é pronunciado.


\begin{center}
\begin{tikzpicture}
  \draw[thick, dashed] (0,0) rectangle (12,6);
\end{tikzpicture}
\end{center}

\noindent\textbf{Escreva 2 palavras com o som da atividade:}

\begin{center}
  \underline{\hspace{3cm}} \quad \underline{\hspace{3cm}}
\end{center}

\noindent\textbf{Desenhe uma situação com a palavra da atividade:}

\begin{center}
\begin{tikzpicture}
  \draw[thick, dashed] (0,0) rectangle (12,4.5);
\end{tikzpicture}
\end{center}

\end{exercicio}

%% ATIVIDADE 7: GUE E GUI: O U MUDO

\plancheta{Atividade de Som 7 --- GUE e GUI: o U mudo}

\subsection{Atividade de Som 7 --- GUE e GUI: o U mudo}
\begin{exercicio}[Atividade de Som 7 --- GUE e GUI: o U mudo]

\metadata{Som07}{Silábico-Alfabético}{EF02LP01, EF02LP04}{D1}{EF02LP04}{CNE/CEB nº 2/2012}

\kumontiming{2}{90}

\noindent\textbf{Nome:} \underline{\hspace{3.2cm}} \quad \textbf{Data:} \underline{\hspace{1.6cm}} \quad \textbf{Nota:} \underline{\hspace{0.9cm}}/10


\noindent\textbf{Comando:} Diga GUERRA sem dizer o U: GUE. Compare com GUITARRA (GUI). Marque as palavras em que o U não é pronunciado.


\begin{center}
\begin{tikzpicture}
  \draw[thick, dashed] (0,0) rectangle (12,6);
\end{tikzpicture}
\end{center}

\noindent\textbf{Escreva 2 palavras com o som da atividade:}

\begin{center}
  \underline{\hspace{3cm}} \quad \underline{\hspace{3cm}}
\end{center}

\noindent\textbf{Desenhe uma situação com a palavra da atividade:}

\begin{center}
\begin{tikzpicture}
  \draw[thick, dashed] (0,0) rectangle (12,4.5);
\end{tikzpicture}
\end{center}

\end{exercicio}

%% ATIVIDADE 8: ENCONTRO BR EM SÍLABAS

\plancheta{Atividade de Som 8 --- Encontro BR em sílabas}

\subsection{Atividade de Som 8 --- Encontro BR em sílabas}
\begin{exercicio}[Atividade de Som 8 --- Encontro BR em sílabas]

\metadata{Som08}{Silábico-Alfabético}{EF02LP01, EF02LP04}{D1}{EF02LP04}{CNE/CEB nº 2/2012}

\kumontiming{2}{90}

\noindent\textbf{Nome:} \underline{\hspace{3.2cm}} \quad \textbf{Data:} \underline{\hspace{1.6cm}} \quad \textbf{Nota:} \underline{\hspace{0.9cm}}/10


\noindent\textbf{Comando:} Diga BRA devagar: B-R-A. Repita mais rápido até virar BRA. Faça o mesmo com BRE, BRI, BRO, BRU.


\begin{center}
\begin{tikzpicture}
  \draw[thick, dashed] (0,0) rectangle (12,6);
\end{tikzpicture}
\end{center}

\noindent\textbf{Escreva 2 palavras com o som da atividade:}

\begin{center}
  \underline{\hspace{3cm}} \quad \underline{\hspace{3cm}}
\end{center}

\noindent\textbf{Desenhe uma situação com a palavra da atividade:}

\begin{center}
\begin{tikzpicture}
  \draw[thick, dashed] (0,0) rectangle (12,4.5);
\end{tikzpicture}
\end{center}

\end{exercicio}

%% ATIVIDADE 9: ENCONTRO CR EM SÍLABAS

\plancheta{Atividade de Som 9 --- Encontro CR em sílabas}

\subsection{Atividade de Som 9 --- Encontro CR em sílabas}
\begin{exercicio}[Atividade de Som 9 --- Encontro CR em sílabas]

\metadata{Som09}{Silábico-Alfabético}{EF02LP01, EF02LP04}{D1}{EF02LP04}{CNE/CEB nº 2/2012}

\kumontiming{2}{90}

\noindent\textbf{Nome:} \underline{\hspace{3.2cm}} \quad \textbf{Data:} \underline{\hspace{1.6cm}} \quad \textbf{Nota:} \underline{\hspace{0.9cm}}/10


\noindent\textbf{Comando:} Diga CRA devagar: C-R-A. Repita mais rápido até virar CRA. O que mudou entre CA e CRA? O R entrou no meio.


\begin{center}
\begin{tikzpicture}
  \draw[thick, dashed] (0,0) rectangle (12,6);
\end{tikzpicture}
\end{center}

\noindent\textbf{Escreva 2 palavras com o som da atividade:}

\begin{center}
  \underline{\hspace{3cm}} \quad \underline{\hspace{3cm}}
\end{center}

\noindent\textbf{Desenhe uma situação com a palavra da atividade:}

\begin{center}
\begin{tikzpicture}
  \draw[thick, dashed] (0,0) rectangle (12,4.5);
\end{tikzpicture}
\end{center}

\end{exercicio}

%% ATIVIDADE 10: ENCONTRO DR EM SÍLABAS

\plancheta{Atividade de Som 10 --- Encontro DR em sílabas}

\subsection{Atividade de Som 10 --- Encontro DR em sílabas}
\begin{exercicio}[Atividade de Som 10 --- Encontro DR em sílabas]

\metadata{Som10}{Silábico-Alfabético}{EF02LP01, EF02LP04}{D1}{EF02LP04}{CNE/CEB nº 2/2012}

\kumontiming{2}{90}

\noindent\textbf{Nome:} \underline{\hspace{3.2cm}} \quad \textbf{Data:} \underline{\hspace{1.6cm}} \quad \textbf{Nota:} \underline{\hspace{0.9cm}}/10


\noindent\textbf{Comando:} Diga DRA devagar: D-R-A. Repita mais rápido. Compare DAMA com DRAMA: o que o R fez?


\begin{center}
\begin{tikzpicture}
  \draw[thick, dashed] (0,0) rectangle (12,6);
\end{tikzpicture}
\end{center}

\noindent\textbf{Escreva 2 palavras com o som da atividade:}

\begin{center}
  \underline{\hspace{3cm}} \quad \underline{\hspace{3cm}}
\end{center}

\noindent\textbf{Desenhe uma situação com a palavra da atividade:}

\begin{center}
\begin{tikzpicture}
  \draw[thick, dashed] (0,0) rectangle (12,4.5);
\end{tikzpicture}
\end{center}

\end{exercicio}

%% ATIVIDADE 11: ENCONTRO FR EM SÍLABAS

\plancheta{Atividade de Som 11 --- Encontro FR em sílabas}

\subsection{Atividade de Som 11 --- Encontro FR em sílabas}
\begin{exercicio}[Atividade de Som 11 --- Encontro FR em sílabas]

\metadata{Som11}{Silábico-Alfabético}{EF02LP01, EF02LP04}{D1}{EF02LP04}{CNE/CEB nº 2/2012}

\kumontiming{2}{90}

\noindent\textbf{Nome:} \underline{\hspace{3.2cm}} \quad \textbf{Data:} \underline{\hspace{1.6cm}} \quad \textbf{Nota:} \underline{\hspace{0.9cm}}/10


\noindent\textbf{Comando:} Diga FRA devagar: F-R-A. Repita mais rápido. Encontre palavras que começam com FRA, FRE, FRI, FRO, FRU.


\begin{center}
\begin{tikzpicture}
  \draw[thick, dashed] (0,0) rectangle (12,6);
\end{tikzpicture}
\end{center}

\noindent\textbf{Escreva 2 palavras com o som da atividade:}

\begin{center}
  \underline{\hspace{3cm}} \quad \underline{\hspace{3cm}}
\end{center}

\noindent\textbf{Desenhe uma situação com a palavra da atividade:}

\begin{center}
\begin{tikzpicture}
  \draw[thick, dashed] (0,0) rectangle (12,4.5);
\end{tikzpicture}
\end{center}

\end{exercicio}

%% ATIVIDADE 12: ENCONTRO GR EM SÍLABAS

\plancheta{Atividade de Som 12 --- Encontro GR em sílabas}

\subsection{Atividade de Som 12 --- Encontro GR em sílabas}
\begin{exercicio}[Atividade de Som 12 --- Encontro GR em sílabas]

\metadata{Som12}{Silábico-Alfabético}{EF02LP01, EF02LP04}{D1}{EF02LP04}{CNE/CEB nº 2/2012}

\kumontiming{2}{90}

\noindent\textbf{Nome:} \underline{\hspace{3.2cm}} \quad \textbf{Data:} \underline{\hspace{1.6cm}} \quad \textbf{Nota:} \underline{\hspace{0.9cm}}/10


\noindent\textbf{Comando:} Diga GRA devagar: G-R-A. Repita mais rápido. O que mudou entre GA e GRA?


\begin{center}
\begin{tikzpicture}
  \draw[thick, dashed] (0,0) rectangle (12,6);
\end{tikzpicture}
\end{center}

\noindent\textbf{Escreva 2 palavras com o som da atividade:}

\begin{center}
  \underline{\hspace{3cm}} \quad \underline{\hspace{3cm}}
\end{center}

\noindent\textbf{Desenhe uma situação com a palavra da atividade:}

\begin{center}
\begin{tikzpicture}
  \draw[thick, dashed] (0,0) rectangle (12,4.5);
\end{tikzpicture}
\end{center}

\end{exercicio}

%% ATIVIDADE 13: ENCONTRO PR EM SÍLABAS

\plancheta{Atividade de Som 13 --- Encontro PR em sílabas}

\subsection{Atividade de Som 13 --- Encontro PR em sílabas}
\begin{exercicio}[Atividade de Som 13 --- Encontro PR em sílabas]

\metadata{Som13}{Silábico-Alfabético}{EF02LP01, EF02LP04}{D1}{EF02LP04}{CNE/CEB nº 2/2012}

\kumontiming{2}{90}

\noindent\textbf{Nome:} \underline{\hspace{3.2cm}} \quad \textbf{Data:} \underline{\hspace{1.6cm}} \quad \textbf{Nota:} \underline{\hspace{0.9cm}}/10


\noindent\textbf{Comando:} Diga PRA devagar: P-R-A. Repita mais rápido. Compare PATO com PRATO: o que o R fez?


\begin{center}
\begin{tikzpicture}
  \draw[thick, dashed] (0,0) rectangle (12,6);
\end{tikzpicture}
\end{center}

\noindent\textbf{Escreva 2 palavras com o som da atividade:}

\begin{center}
  \underline{\hspace{3cm}} \quad \underline{\hspace{3cm}}
\end{center}

\noindent\textbf{Desenhe uma situação com a palavra da atividade:}

\begin{center}
\begin{tikzpicture}
  \draw[thick, dashed] (0,0) rectangle (12,4.5);
\end{tikzpicture}
\end{center}

\end{exercicio}

%% ATIVIDADE 14: ENCONTRO TR EM SÍLABAS

\plancheta{Atividade de Som 14 --- Encontro TR em sílabas}

\subsection{Atividade de Som 14 --- Encontro TR em sílabas}
\begin{exercicio}[Atividade de Som 14 --- Encontro TR em sílabas]

\metadata{Som14}{Silábico-Alfabético}{EF02LP01, EF02LP04}{D1}{EF02LP04}{CNE/CEB nº 2/2012}

\kumontiming{2}{90}

\noindent\textbf{Nome:} \underline{\hspace{3.2cm}} \quad \textbf{Data:} \underline{\hspace{1.6cm}} \quad \textbf{Nota:} \underline{\hspace{0.9cm}}/10


\noindent\textbf{Comando:} Diga TRA devagar: T-R-A. Repita mais rápido. Compare TAPA com TRAPA: existe TRAPA?


\begin{center}
\begin{tikzpicture}
  \draw[thick, dashed] (0,0) rectangle (12,6);
\end{tikzpicture}
\end{center}

\noindent\textbf{Escreva 2 palavras com o som da atividade:}

\begin{center}
  \underline{\hspace{3cm}} \quad \underline{\hspace{3cm}}
\end{center}

\noindent\textbf{Desenhe uma situação com a palavra da atividade:}

\begin{center}
\begin{tikzpicture}
  \draw[thick, dashed] (0,0) rectangle (12,4.5);
\end{tikzpicture}
\end{center}

\end{exercicio}

%% ATIVIDADE 15: ENCONTRO VR EM SÍLABAS

\plancheta{Atividade de Som 15 --- Encontro VR em sílabas}

\subsection{Atividade de Som 15 --- Encontro VR em sílabas}
\begin{exercicio}[Atividade de Som 15 --- Encontro VR em sílabas]

\metadata{Som15}{Silábico-Alfabético}{EF02LP01, EF02LP04}{D1}{EF02LP04}{CNE/CEB nº 2/2012}

\kumontiming{2}{90}

\noindent\textbf{Nome:} \underline{\hspace{3.2cm}} \quad \textbf{Data:} \underline{\hspace{1.6cm}} \quad \textbf{Nota:} \underline{\hspace{0.9cm}}/10


\noindent\textbf{Comando:} Diga VRA devagar: V-R-A. Repita mais rápido. Encontre LIVRO e PALAVRA: onde está o VR?


\begin{center}
\begin{tikzpicture}
  \draw[thick, dashed] (0,0) rectangle (12,6);
\end{tikzpicture}
\end{center}

\noindent\textbf{Escreva 2 palavras com o som da atividade:}

\begin{center}
  \underline{\hspace{3cm}} \quad \underline{\hspace{3cm}}
\end{center}

\noindent\textbf{Desenhe uma situação com a palavra da atividade:}

\begin{center}
\begin{tikzpicture}
  \draw[thick, dashed] (0,0) rectangle (12,4.5);
\end{tikzpicture}
\end{center}

\end{exercicio}

%% ATIVIDADE 16: ENCONTRO BL EM SÍLABAS

\plancheta{Atividade de Som 16 --- Encontro BL em sílabas}

\subsection{Atividade de Som 16 --- Encontro BL em sílabas}
\begin{exercicio}[Atividade de Som 16 --- Encontro BL em sílabas]

\metadata{Som16}{Silábico-Alfabético}{EF02LP01, EF02LP04}{D1}{EF02LP04}{CNE/CEB nº 2/2012}

\kumontiming{2}{90}

\noindent\textbf{Nome:} \underline{\hspace{3.2cm}} \quad \textbf{Data:} \underline{\hspace{1.6cm}} \quad \textbf{Nota:} \underline{\hspace{0.9cm}}/10


\noindent\textbf{Comando:} Diga BLA devagar: B-L-A. Repita mais rápido. O L vem logo depois do B, sem vogal no meio.


\begin{center}
\begin{tikzpicture}
  \draw[thick, dashed] (0,0) rectangle (12,6);
\end{tikzpicture}
\end{center}

\noindent\textbf{Escreva 2 palavras com o som da atividade:}

\begin{center}
  \underline{\hspace{3cm}} \quad \underline{\hspace{3cm}}
\end{center}

\noindent\textbf{Desenhe uma situação com a palavra da atividade:}

\begin{center}
\begin{tikzpicture}
  \draw[thick, dashed] (0,0) rectangle (12,4.5);
\end{tikzpicture}
\end{center}

\end{exercicio}

%% ATIVIDADE 17: ENCONTRO CL EM SÍLABAS

\plancheta{Atividade de Som 17 --- Encontro CL em sílabas}

\subsection{Atividade de Som 17 --- Encontro CL em sílabas}
\begin{exercicio}[Atividade de Som 17 --- Encontro CL em sílabas]

\metadata{Som17}{Silábico-Alfabético}{EF02LP01, EF02LP04}{D1}{EF02LP04}{CNE/CEB nº 2/2012}

\kumontiming{2}{90}

\noindent\textbf{Nome:} \underline{\hspace{3.2cm}} \quad \textbf{Data:} \underline{\hspace{1.6cm}} \quad \textbf{Nota:} \underline{\hspace{0.9cm}}/10


\noindent\textbf{Comando:} Diga CLA devagar: C-L-A. Repita mais rápido. Compare CASA com CLASSE: o L mudou o som?


\begin{center}
\begin{tikzpicture}
  \draw[thick, dashed] (0,0) rectangle (12,6);
\end{tikzpicture}
\end{center}

\noindent\textbf{Escreva 2 palavras com o som da atividade:}

\begin{center}
  \underline{\hspace{3cm}} \quad \underline{\hspace{3cm}}
\end{center}

\noindent\textbf{Desenhe uma situação com a palavra da atividade:}

\begin{center}
\begin{tikzpicture}
  \draw[thick, dashed] (0,0) rectangle (12,4.5);
\end{tikzpicture}
\end{center}

\end{exercicio}

%% ATIVIDADE 18: ENCONTRO FL EM SÍLABAS

\plancheta{Atividade de Som 18 --- Encontro FL em sílabas}

\subsection{Atividade de Som 18 --- Encontro FL em sílabas}
\begin{exercicio}[Atividade de Som 18 --- Encontro FL em sílabas]

\metadata{Som18}{Silábico-Alfabético}{EF02LP01, EF02LP04}{D1}{EF02LP04}{CNE/CEB nº 2/2012}

\kumontiming{2}{90}

\noindent\textbf{Nome:} \underline{\hspace{3.2cm}} \quad \textbf{Data:} \underline{\hspace{1.6cm}} \quad \textbf{Nota:} \underline{\hspace{0.9cm}}/10


\noindent\textbf{Comando:} Diga FLA devagar: F-L-A. Repita mais rápido. FLOR e FLAUTA: onde está o FL?


\begin{center}
\begin{tikzpicture}
  \draw[thick, dashed] (0,0) rectangle (12,6);
\end{tikzpicture}
\end{center}

\noindent\textbf{Escreva 2 palavras com o som da atividade:}

\begin{center}
  \underline{\hspace{3cm}} \quad \underline{\hspace{3cm}}
\end{center}

\noindent\textbf{Desenhe uma situação com a palavra da atividade:}

\begin{center}
\begin{tikzpicture}
  \draw[thick, dashed] (0,0) rectangle (12,4.5);
\end{tikzpicture}
\end{center}

\end{exercicio}

%% ATIVIDADE 19: ENCONTRO GL EM SÍLABAS

\plancheta{Atividade de Som 19 --- Encontro GL em sílabas}

\subsection{Atividade de Som 19 --- Encontro GL em sílabas}
\begin{exercicio}[Atividade de Som 19 --- Encontro GL em sílabas]

\metadata{Som19}{Silábico-Alfabético}{EF02LP01, EF02LP04}{D1}{EF02LP04}{CNE/CEB nº 2/2012}

\kumontiming{2}{90}

\noindent\textbf{Nome:} \underline{\hspace{3.2cm}} \quad \textbf{Data:} \underline{\hspace{1.6cm}} \quad \textbf{Nota:} \underline{\hspace{0.9cm}}/10


\noindent\textbf{Comando:} Diga GLA devagar: G-L-A. Repita mais rápido. GLOBO e GLÓRIA: onde está o GL?


\begin{center}
\begin{tikzpicture}
  \draw[thick, dashed] (0,0) rectangle (12,6);
\end{tikzpicture}
\end{center}

\noindent\textbf{Escreva 2 palavras com o som da atividade:}

\begin{center}
  \underline{\hspace{3cm}} \quad \underline{\hspace{3cm}}
\end{center}

\noindent\textbf{Desenhe uma situação com a palavra da atividade:}

\begin{center}
\begin{tikzpicture}
  \draw[thick, dashed] (0,0) rectangle (12,4.5);
\end{tikzpicture}
\end{center}

\end{exercicio}

%% ATIVIDADE 20: ENCONTRO PL EM SÍLABAS

\plancheta{Atividade de Som 20 --- Encontro PL em sílabas}

\subsection{Atividade de Som 20 --- Encontro PL em sílabas}
\begin{exercicio}[Atividade de Som 20 --- Encontro PL em sílabas]

\metadata{Som20}{Silábico-Alfabético}{EF02LP01, EF02LP04}{D1}{EF02LP04}{CNE/CEB nº 2/2012}

\kumontiming{2}{90}

\noindent\textbf{Nome:} \underline{\hspace{3.2cm}} \quad \textbf{Data:} \underline{\hspace{1.6cm}} \quad \textbf{Nota:} \underline{\hspace{0.9cm}}/10


\noindent\textbf{Comando:} Diga PLA devagar: P-L-A. Repita mais rápido. PLACA e PLANTA: onde está o PL?


\begin{center}
\begin{tikzpicture}
  \draw[thick, dashed] (0,0) rectangle (12,6);
\end{tikzpicture}
\end{center}

\noindent\textbf{Escreva 2 palavras com o som da atividade:}

\begin{center}
  \underline{\hspace{3cm}} \quad \underline{\hspace{3cm}}
\end{center}

\noindent\textbf{Desenhe uma situação com a palavra da atividade:}

\begin{center}
\begin{tikzpicture}
  \draw[thick, dashed] (0,0) rectangle (12,4.5);
\end{tikzpicture}
\end{center}

\end{exercicio}

%% ATIVIDADE 21: R FORTE E R FRACO

\plancheta{Atividade de Som 21 --- R forte e R fraco}

\subsection{Atividade de Som 21 --- R forte e R fraco}
\begin{exercicio}[Atividade de Som 21 --- R forte e R fraco]

\metadata{Som21}{Silábico-Alfabético}{EF02LP01, EF02LP04}{D1}{EF02LP04}{CNE/CEB nº 2/2012}

\kumontiming{2}{90}

\noindent\textbf{Nome:} \underline{\hspace{3.2cm}} \quad \textbf{Data:} \underline{\hspace{1.6cm}} \quad \textbf{Nota:} \underline{\hspace{0.9cm}}/10


\noindent\textbf{Comando:} Compare CARA (R fraco) com CARRO (R forte). Diga as palavras e circule onde o R vibra mais.


\begin{center}
\begin{tikzpicture}
  \draw[thick, dashed] (0,0) rectangle (12,6);
\end{tikzpicture}
\end{center}

\noindent\textbf{Escreva 2 palavras com o som da atividade:}

\begin{center}
  \underline{\hspace{3cm}} \quad \underline{\hspace{3cm}}
\end{center}

\noindent\textbf{Desenhe uma situação com a palavra da atividade:}

\begin{center}
\begin{tikzpicture}
  \draw[thick, dashed] (0,0) rectangle (12,4.5);
\end{tikzpicture}
\end{center}

\end{exercicio}

%% ATIVIDADE 22: SS ENTRE VOGAIS

\plancheta{Atividade de Som 22 --- SS entre vogais}

\subsection{Atividade de Som 22 --- SS entre vogais}
\begin{exercicio}[Atividade de Som 22 --- SS entre vogais]

\metadata{Som22}{Silábico-Alfabético}{EF02LP01, EF02LP04}{D1}{EF02LP04}{CNE/CEB nº 2/2012}

\kumontiming{2}{90}

\noindent\textbf{Nome:} \underline{\hspace{3.2cm}} \quad \textbf{Data:} \underline{\hspace{1.6cm}} \quad \textbf{Nota:} \underline{\hspace{0.9cm}}/10


\noindent\textbf{Comando:} Compare CASA (S com som de Z) com PASSA (SS com som de S). O SS sempre mantém o som de cobra.


\begin{center}
\begin{tikzpicture}
  \draw[thick, dashed] (0,0) rectangle (12,6);
\end{tikzpicture}
\end{center}

\noindent\textbf{Escreva 2 palavras com o som da atividade:}

\begin{center}
  \underline{\hspace{3cm}} \quad \underline{\hspace{3cm}}
\end{center}

\noindent\textbf{Desenhe uma situação com a palavra da atividade:}

\begin{center}
\begin{tikzpicture}
  \draw[thick, dashed] (0,0) rectangle (12,4.5);
\end{tikzpicture}
\end{center}

\end{exercicio}

%% ATIVIDADE 23: S ENTRE VOGAIS COM SOM DE Z

\plancheta{Atividade de Som 23 --- S entre vogais com som de Z}

\subsection{Atividade de Som 23 --- S entre vogais com som de Z}
\begin{exercicio}[Atividade de Som 23 --- S entre vogais com som de Z]

\metadata{Som23}{Silábico-Alfabético}{EF02LP01, EF02LP04}{D1}{EF02LP04}{CNE/CEB nº 2/2012}

\kumontiming{2}{90}

\noindent\textbf{Nome:} \underline{\hspace{3.2cm}} \quad \textbf{Data:} \underline{\hspace{1.6cm}} \quad \textbf{Nota:} \underline{\hspace{0.9cm}}/10


\noindent\textbf{Comando:} Diga CASA, MESA, ASA. Perceba: o S ficou com som de Z. Esta é uma das regras mais importantes do ano.


\begin{center}
\begin{tikzpicture}
  \draw[thick, dashed] (0,0) rectangle (12,6);
\end{tikzpicture}
\end{center}

\noindent\textbf{Escreva 2 palavras com o som da atividade:}

\begin{center}
  \underline{\hspace{3cm}} \quad \underline{\hspace{3cm}}
\end{center}

\noindent\textbf{Desenhe uma situação com a palavra da atividade:}

\begin{center}
\begin{tikzpicture}
  \draw[thick, dashed] (0,0) rectangle (12,4.5);
\end{tikzpicture}
\end{center}

\end{exercicio}

%% ATIVIDADE 24: Z EM POSIÇÃO INICIAL

\plancheta{Atividade de Som 24 --- Z em posição inicial}

\subsection{Atividade de Som 24 --- Z em posição inicial}
\begin{exercicio}[Atividade de Som 24 --- Z em posição inicial]

\metadata{Som24}{Silábico-Alfabético}{EF02LP01, EF02LP04}{D1}{EF02LP04}{CNE/CEB nº 2/2012}

\kumontiming{2}{90}

\noindent\textbf{Nome:} \underline{\hspace{3.2cm}} \quad \textbf{Data:} \underline{\hspace{1.6cm}} \quad \textbf{Nota:} \underline{\hspace{0.9cm}}/10


\noindent\textbf{Comando:} Diga ZERO, ZEBRA, ZÍPER. O Z é uma letra que 'liga' a voz. Escreva outras palavras com Z.


\begin{center}
\begin{tikzpicture}
  \draw[thick, dashed] (0,0) rectangle (12,6);
\end{tikzpicture}
\end{center}

\noindent\textbf{Escreva 2 palavras com o som da atividade:}

\begin{center}
  \underline{\hspace{3cm}} \quad \underline{\hspace{3cm}}
\end{center}

\noindent\textbf{Desenhe uma situação com a palavra da atividade:}

\begin{center}
\begin{tikzpicture}
  \draw[thick, dashed] (0,0) rectangle (12,4.5);
\end{tikzpicture}
\end{center}

\end{exercicio}

%% ATIVIDADE 25: Ç EM PALAVRAS

\plancheta{Atividade de Som 25 --- Ç em palavras}

\subsection{Atividade de Som 25 --- Ç em palavras}
\begin{exercicio}[Atividade de Som 25 --- Ç em palavras]

\metadata{Som25}{Silábico-Alfabético}{EF02LP01, EF02LP04}{D1}{EF02LP04}{CNE/CEB nº 2/2012}

\kumontiming{2}{90}

\noindent\textbf{Nome:} \underline{\hspace{3.2cm}} \quad \textbf{Data:} \underline{\hspace{1.6cm}} \quad \textbf{Nota:} \underline{\hspace{0.9cm}}/10


\noindent\textbf{Comando:} Diga CORAÇÃO, AÇÚCAR, LAÇO. O Ç tem som de S. Escreva palavras com ÇA, ÇO, ÇU.


\begin{center}
\begin{tikzpicture}
  \draw[thick, dashed] (0,0) rectangle (12,6);
\end{tikzpicture}
\end{center}

\noindent\textbf{Escreva 2 palavras com o som da atividade:}

\begin{center}
  \underline{\hspace{3cm}} \quad \underline{\hspace{3cm}}
\end{center}

\noindent\textbf{Desenhe uma situação com a palavra da atividade:}

\begin{center}
\begin{tikzpicture}
  \draw[thick, dashed] (0,0) rectangle (12,4.5);
\end{tikzpicture}
\end{center}

\end{exercicio}

%% ATIVIDADE 26: AM NASAL NO FINAL

\plancheta{Atividade de Som 26 --- AM nasal no final}

\subsection{Atividade de Som 26 --- AM nasal no final}
\begin{exercicio}[Atividade de Som 26 --- AM nasal no final]

\metadata{Som26}{Silábico-Alfabético}{EF02LP01, EF02LP04}{D1}{EF02LP04}{CNE/CEB nº 2/2012}

\kumontiming{2}{90}

\noindent\textbf{Nome:} \underline{\hspace{3.2cm}} \quad \textbf{Data:} \underline{\hspace{1.6cm}} \quad \textbf{Nota:} \underline{\hspace{0.9cm}}/10


\noindent\textbf{Comando:} Diga CANTAM, FALAM, BRINCAM. O AM final é a desinência de 'eles'. Sinta o nariz vibrar.


\begin{center}
\begin{tikzpicture}
  \draw[thick, dashed] (0,0) rectangle (12,6);
\end{tikzpicture}
\end{center}

\noindent\textbf{Escreva 2 palavras com o som da atividade:}

\begin{center}
  \underline{\hspace{3cm}} \quad \underline{\hspace{3cm}}
\end{center}

\noindent\textbf{Desenhe uma situação com a palavra da atividade:}

\begin{center}
\begin{tikzpicture}
  \draw[thick, dashed] (0,0) rectangle (12,4.5);
\end{tikzpicture}
\end{center}

\end{exercicio}

%% ATIVIDADE 27: AN NASAL NO MEIO

\plancheta{Atividade de Som 27 --- AN nasal no meio}

\subsection{Atividade de Som 27 --- AN nasal no meio}
\begin{exercicio}[Atividade de Som 27 --- AN nasal no meio]

\metadata{Som27}{Silábico-Alfabético}{EF02LP01, EF02LP04}{D1}{EF02LP04}{CNE/CEB nº 2/2012}

\kumontiming{2}{90}

\noindent\textbf{Nome:} \underline{\hspace{3.2cm}} \quad \textbf{Data:} \underline{\hspace{1.6cm}} \quad \textbf{Nota:} \underline{\hspace{0.9cm}}/10


\noindent\textbf{Comando:} Diga BANCO, CAMPO, DANÇA. O AN é uma sílaba nasal no meio da palavra.


\begin{center}
\begin{tikzpicture}
  \draw[thick, dashed] (0,0) rectangle (12,6);
\end{tikzpicture}
\end{center}

\noindent\textbf{Escreva 2 palavras com o som da atividade:}

\begin{center}
  \underline{\hspace{3cm}} \quad \underline{\hspace{3cm}}
\end{center}

\noindent\textbf{Desenhe uma situação com a palavra da atividade:}

\begin{center}
\begin{tikzpicture}
  \draw[thick, dashed] (0,0) rectangle (12,4.5);
\end{tikzpicture}
\end{center}

\end{exercicio}

%% ATIVIDADE 28: X COM SOM DE CH

\plancheta{Atividade de Som 28 --- X com som de CH}

\subsection{Atividade de Som 28 --- X com som de CH}
\begin{exercicio}[Atividade de Som 28 --- X com som de CH]

\metadata{Som28}{Silábico-Alfabético}{EF02LP01, EF02LP04}{D1}{EF02LP04}{CNE/CEB nº 2/2012}

\kumontiming{2}{90}

\noindent\textbf{Nome:} \underline{\hspace{3.2cm}} \quad \textbf{Data:} \underline{\hspace{1.6cm}} \quad \textbf{Nota:} \underline{\hspace{0.9cm}}/10


\noindent\textbf{Comando:} Diga XADREZ, PEIXE, CAIXA. O X 'veste' o som de CH. Leia e desenhe uma caixa.


\begin{center}
\begin{tikzpicture}
  \draw[thick, dashed] (0,0) rectangle (12,6);
\end{tikzpicture}
\end{center}

\noindent\textbf{Escreva 2 palavras com o som da atividade:}

\begin{center}
  \underline{\hspace{3cm}} \quad \underline{\hspace{3cm}}
\end{center}

\noindent\textbf{Desenhe uma situação com a palavra da atividade:}

\begin{center}
\begin{tikzpicture}
  \draw[thick, dashed] (0,0) rectangle (12,4.5);
\end{tikzpicture}
\end{center}

\end{exercicio}

%% ATIVIDADE 29: X COM SOM DE Z

\plancheta{Atividade de Som 29 --- X com som de Z}

\subsection{Atividade de Som 29 --- X com som de Z}
\begin{exercicio}[Atividade de Som 29 --- X com som de Z]

\metadata{Som29}{Silábico-Alfabético}{EF02LP01, EF02LP04}{D1}{EF02LP04}{CNE/CEB nº 2/2012}

\kumontiming{2}{90}

\noindent\textbf{Nome:} \underline{\hspace{3.2cm}} \quad \textbf{Data:} \underline{\hspace{1.6cm}} \quad \textbf{Nota:} \underline{\hspace{0.9cm}}/10


\noindent\textbf{Comando:} Diga EXAME, EXEMPLO. O X 'veste' o som de Z em palavras com EXA, EXE.


\begin{center}
\begin{tikzpicture}
  \draw[thick, dashed] (0,0) rectangle (12,6);
\end{tikzpicture}
\end{center}

\noindent\textbf{Escreva 2 palavras com o som da atividade:}

\begin{center}
  \underline{\hspace{3cm}} \quad \underline{\hspace{3cm}}
\end{center}

\noindent\textbf{Desenhe uma situação com a palavra da atividade:}

\begin{center}
\begin{tikzpicture}
  \draw[thick, dashed] (0,0) rectangle (12,4.5);
\end{tikzpicture}
\end{center}

\end{exercicio}

%% ATIVIDADE 30: X COM SOM DE KS

\plancheta{Atividade de Som 30 --- X com som de KS}

\subsection{Atividade de Som 30 --- X com som de KS}
\begin{exercicio}[Atividade de Som 30 --- X com som de KS]

\metadata{Som30}{Silábico-Alfabético}{EF02LP01, EF02LP04}{D1}{EF02LP04}{CNE/CEB nº 2/2012}

\kumontiming{2}{90}

\noindent\textbf{Nome:} \underline{\hspace{3.2cm}} \quad \textbf{Data:} \underline{\hspace{1.6cm}} \quad \textbf{Nota:} \underline{\hspace{0.9cm}}/10


\noindent\textbf{Comando:} Diga TÁXI, SAXOFONE. O X 'veste' os sons de K e S juntos. Como pronunciar FIXO?


\begin{center}
\begin{tikzpicture}
  \draw[thick, dashed] (0,0) rectangle (12,6);
\end{tikzpicture}
\end{center}

\noindent\textbf{Escreva 2 palavras com o som da atividade:}

\begin{center}
  \underline{\hspace{3cm}} \quad \underline{\hspace{3cm}}
\end{center}

\noindent\textbf{Desenhe uma situação com a palavra da atividade:}

\begin{center}
\begin{tikzpicture}
  \draw[thick, dashed] (0,0) rectangle (12,4.5);
\end{tikzpicture}
\end{center}

\end{exercicio}

%% ATIVIDADE 31: X COM SOM DE SS

\plancheta{Atividade de Som 31 --- X com som de SS}

\subsection{Atividade de Som 31 --- X com som de SS}
\begin{exercicio}[Atividade de Som 31 --- X com som de SS]

\metadata{Som31}{Silábico-Alfabético}{EF02LP01, EF02LP04}{D1}{EF02LP04}{CNE/CEB nº 2/2012}

\kumontiming{2}{90}

\noindent\textbf{Nome:} \underline{\hspace{3.2cm}} \quad \textbf{Data:} \underline{\hspace{1.6cm}} \quad \textbf{Nota:} \underline{\hspace{0.9cm}}/10


\noindent\textbf{Comando:} Diga TEXTO, EXTRA. O X 'veste' o som de S. Leia em voz alta e marque a pronúncia.


\begin{center}
\begin{tikzpicture}
  \draw[thick, dashed] (0,0) rectangle (12,6);
\end{tikzpicture}
\end{center}

\noindent\textbf{Escreva 2 palavras com o som da atividade:}

\begin{center}
  \underline{\hspace{3cm}} \quad \underline{\hspace{3cm}}
\end{center}

\noindent\textbf{Desenhe uma situação com a palavra da atividade:}

\begin{center}
\begin{tikzpicture}
  \draw[thick, dashed] (0,0) rectangle (12,4.5);
\end{tikzpicture}
\end{center}

\end{exercicio}

%% ATIVIDADE 32: RIMA E ALITERAÇÃO

\plancheta{Atividade de Som 32 --- Rima e aliteração}

\subsection{Atividade de Som 32 --- Rima e aliteração}
\begin{exercicio}[Atividade de Som 32 --- Rima e aliteração]

\metadata{Som32}{Silábico-Alfabético}{EF02LP01, EF02LP04}{D1}{EF02LP04}{CNE/CEB nº 2/2012}

\kumontiming{2}{90}

\noindent\textbf{Nome:} \underline{\hspace{3.2cm}} \quad \textbf{Data:} \underline{\hspace{1.6cm}} \quad \textbf{Nota:} \underline{\hspace{0.9cm}}/10


\noindent\textbf{Comando:} Ouça: CHAVE, CHUVA, CHÃO. Todas começam com o mesmo som. Crie outra série com LH.


\begin{center}
\begin{tikzpicture}
  \draw[thick, dashed] (0,0) rectangle (12,6);
\end{tikzpicture}
\end{center}

\noindent\textbf{Escreva 2 palavras com o som da atividade:}

\begin{center}
  \underline{\hspace{3cm}} \quad \underline{\hspace{3cm}}
\end{center}

\noindent\textbf{Desenhe uma situação com a palavra da atividade:}

\begin{center}
\begin{tikzpicture}
  \draw[thick, dashed] (0,0) rectangle (12,4.5);
\end{tikzpicture}
\end{center}

\end{exercicio}

%% ATIVIDADE 33: SEGMENTAÇÃO SILÁBICA AVANÇADA

\plancheta{Atividade de Som 33 --- Segmentação silábica avançada}

\subsection{Atividade de Som 33 --- Segmentação silábica avançada}
\begin{exercicio}[Atividade de Som 33 --- Segmentação silábica avançada]

\metadata{Som33}{Silábico-Alfabético}{EF02LP01, EF02LP04}{D1}{EF02LP04}{CNE/CEB nº 2/2012}

\kumontiming{2}{90}

\noindent\textbf{Nome:} \underline{\hspace{3.2cm}} \quad \textbf{Data:} \underline{\hspace{1.6cm}} \quad \textbf{Nota:} \underline{\hspace{0.9cm}}/10


\noindent\textbf{Comando:} Diga BIBLIOTECA: BI-BLIO-TE-CA. Quantas sílabas? Agora escreva a palavra e sublinhe o encontro consonantal.


\begin{center}
\begin{tikzpicture}
  \draw[thick, dashed] (0,0) rectangle (12,6);
\end{tikzpicture}
\end{center}

\noindent\textbf{Escreva 2 palavras com o som da atividade:}

\begin{center}
  \underline{\hspace{3cm}} \quad \underline{\hspace{3cm}}
\end{center}

\noindent\textbf{Desenhe uma situação com a palavra da atividade:}

\begin{center}
\begin{tikzpicture}
  \draw[thick, dashed] (0,0) rectangle (12,4.5);
\end{tikzpicture}
\end{center}

\end{exercicio}

%% ATIVIDADE 34: POSIÇÃO DO SOM NA PALAVRA

\plancheta{Atividade de Som 34 --- Posição do som na palavra}

\subsection{Atividade de Som 34 --- Posição do som na palavra}
\begin{exercicio}[Atividade de Som 34 --- Posição do som na palavra]

\metadata{Som34}{Silábico-Alfabético}{EF02LP01, EF02LP04}{D1}{EF02LP04}{CNE/CEB nº 2/2012}

\kumontiming{2}{90}

\noindent\textbf{Nome:} \underline{\hspace{3.2cm}} \quad \textbf{Data:} \underline{\hspace{1.6cm}} \quad \textbf{Nota:} \underline{\hspace{0.9cm}}/10


\noindent\textbf{Comando:} Diga GUITARRA: onde está o som de GUE? No início (GUI), no meio (TA) ou no fim (RRA)?


\begin{center}
\begin{tikzpicture}
  \draw[thick, dashed] (0,0) rectangle (12,6);
\end{tikzpicture}
\end{center}

\noindent\textbf{Escreva 2 palavras com o som da atividade:}

\begin{center}
  \underline{\hspace{3cm}} \quad \underline{\hspace{3cm}}
\end{center}

\noindent\textbf{Desenhe uma situação com a palavra da atividade:}

\begin{center}
\begin{tikzpicture}
  \draw[thick, dashed] (0,0) rectangle (12,4.5);
\end{tikzpicture}
\end{center}

\end{exercicio}

%% ATIVIDADE 35: COMPLETAR COM A GRAFIA CERTA

\plancheta{Atividade de Som 35 --- Completar com a grafia certa}

\subsection{Atividade de Som 35 --- Completar com a grafia certa}
\begin{exercicio}[Atividade de Som 35 --- Completar com a grafia certa]

\metadata{Som35}{Silábico-Alfabético}{EF02LP01, EF02LP04}{D1}{EF02LP04}{CNE/CEB nº 2/2012}

\kumontiming{2}{90}

\noindent\textbf{Nome:} \underline{\hspace{3.2cm}} \quad \textbf{Data:} \underline{\hspace{1.6cm}} \quad \textbf{Nota:} \underline{\hspace{0.9cm}}/10


\noindent\textbf{Comando:} Complete: \_AVE (CHAVE), \_APÉU (CHAPÉU), \_UVA (CHUVA). Agora invente você mais três.


\begin{center}
\begin{tikzpicture}
  \draw[thick, dashed] (0,0) rectangle (12,6);
\end{tikzpicture}
\end{center}

\noindent\textbf{Escreva 2 palavras com o som da atividade:}

\begin{center}
  \underline{\hspace{3cm}} \quad \underline{\hspace{3cm}}
\end{center}

\noindent\textbf{Desenhe uma situação com a palavra da atividade:}

\begin{center}
\begin{tikzpicture}
  \draw[thick, dashed] (0,0) rectangle (12,4.5);
\end{tikzpicture}
\end{center}

\end{exercicio}

%% ATIVIDADE 36: DITADO DE GRAFIAS DO ANO

\plancheta{Atividade de Som 36 --- Ditado de grafias do ano}

\subsection{Atividade de Som 36 --- Ditado de grafias do ano}
\begin{exercicio}[Atividade de Som 36 --- Ditado de grafias do ano]

\metadata{Som36}{Silábico-Alfabético}{EF02LP01, EF02LP04}{D1}{EF02LP04}{CNE/CEB nº 2/2012}

\kumontiming{2}{90}

\noindent\textbf{Nome:} \underline{\hspace{3.2cm}} \quad \textbf{Data:} \underline{\hspace{1.6cm}} \quad \textbf{Nota:} \underline{\hspace{0.9cm}}/10


\noindent\textbf{Comando:} Faça o ditado: queijo, coelho, ninho, carro, casa, coração. Corrija com o aluno e registre os acertos.


\begin{center}
\begin{tikzpicture}
  \draw[thick, dashed] (0,0) rectangle (12,6);
\end{tikzpicture}
\end{center}

\noindent\textbf{Escreva 2 palavras com o som da atividade:}

\begin{center}
  \underline{\hspace{3cm}} \quad \underline{\hspace{3cm}}
\end{center}

\noindent\textbf{Desenhe uma situação com a palavra da atividade:}

\begin{center}
\begin{tikzpicture}
  \draw[thick, dashed] (0,0) rectangle (12,4.5);
\end{tikzpicture}
\end{center}

\end{exercicio}

%% ATIVIDADE 37: SUBTRAÇÃO DE FONEMA EM ENCONTROS

\plancheta{Atividade de Som 37 --- Subtração de fonema em encontros}

\subsection{Atividade de Som 37 --- Subtração de fonema em encontros}
\begin{exercicio}[Atividade de Som 37 --- Subtração de fonema em encontros]

\metadata{Som37}{Silábico-Alfabético}{EF02LP01, EF02LP04}{D1}{EF02LP04}{CNE/CEB nº 2/2012}

\kumontiming{2}{90}

\noindent\textbf{Nome:} \underline{\hspace{3.2cm}} \quad \textbf{Data:} \underline{\hspace{1.6cm}} \quad \textbf{Nota:} \underline{\hspace{0.9cm}}/10


\noindent\textbf{Comando:} Diga PRATO e retire o R: PATO. Diga FRIO e retire o R: FIO. Diga BLUSA e retire o L: BUSA? Teste e perceba: nem sempre a palavra continua existindo.


\begin{center}
\begin{tikzpicture}
  \draw[thick, dashed] (0,0) rectangle (12,6);
\end{tikzpicture}
\end{center}

\noindent\textbf{Escreva 2 palavras com o som da atividade:}

\begin{center}
  \underline{\hspace{3cm}} \quad \underline{\hspace{3cm}}
\end{center}

\noindent\textbf{Desenhe uma situação com a palavra da atividade:}

\begin{center}
\begin{tikzpicture}
  \draw[thick, dashed] (0,0) rectangle (12,4.5);
\end{tikzpicture}
\end{center}

\end{exercicio}

%% ATIVIDADE 38: RIMAS AVANÇADAS

\plancheta{Atividade de Som 38 --- Rimas avançadas}

\subsection{Atividade de Som 38 --- Rimas avançadas}
\begin{exercicio}[Atividade de Som 38 --- Rimas avançadas]

\metadata{Som38}{Silábico-Alfabético}{EF02LP01, EF02LP04}{D1}{EF02LP04}{CNE/CEB nº 2/2012}

\kumontiming{2}{90}

\noindent\textbf{Nome:} \underline{\hspace{3.2cm}} \quad \textbf{Data:} \underline{\hspace{1.6cm}} \quad \textbf{Nota:} \underline{\hspace{0.9cm}}/10


\noindent\textbf{Comando:} Encontre rimas: CHUVA (LUVA), PRATO (GATO), TRILHO (MILHO), FLOR (AMOR), CAIXA (FAIXA). Crie uma quadrinha com uma das rimas.


\begin{center}
\begin{tikzpicture}
  \draw[thick, dashed] (0,0) rectangle (12,6);
\end{tikzpicture}
\end{center}

\noindent\textbf{Escreva 2 palavras com o som da atividade:}

\begin{center}
  \underline{\hspace{3cm}} \quad \underline{\hspace{3cm}}
\end{center}

\noindent\textbf{Desenhe uma situação com a palavra da atividade:}

\begin{center}
\begin{tikzpicture}
  \draw[thick, dashed] (0,0) rectangle (12,4.5);
\end{tikzpicture}
\end{center}

\end{exercicio}

%% ATIVIDADE 39: ALITERAÇÃO EM CADEIA

\plancheta{Atividade de Som 39 --- Aliteração em cadeia}

\subsection{Atividade de Som 39 --- Aliteração em cadeia}
\begin{exercicio}[Atividade de Som 39 --- Aliteração em cadeia]

\metadata{Som39}{Silábico-Alfabético}{EF02LP01, EF02LP04}{D1}{EF02LP04}{CNE/CEB nº 2/2012}

\kumontiming{2}{90}

\noindent\textbf{Nome:} \underline{\hspace{3.2cm}} \quad \textbf{Data:} \underline{\hspace{1.6cm}} \quad \textbf{Nota:} \underline{\hspace{0.9cm}}/10


\noindent\textbf{Comando:} Crie uma frase em que todas as palavras comecem com o mesmo som: 'O coelho comeu cenoura no campo' (c). Depois com CH: 'A chuva chovera chão'.


\begin{center}
\begin{tikzpicture}
  \draw[thick, dashed] (0,0) rectangle (12,6);
\end{tikzpicture}
\end{center}

\noindent\textbf{Escreva 2 palavras com o som da atividade:}

\begin{center}
  \underline{\hspace{3cm}} \quad \underline{\hspace{3cm}}
\end{center}

\noindent\textbf{Desenhe uma situação com a palavra da atividade:}

\begin{center}
\begin{tikzpicture}
  \draw[thick, dashed] (0,0) rectangle (12,4.5);
\end{tikzpicture}
\end{center}

\end{exercicio}

%% ATIVIDADE 40: CONTAGEM DE FONEMAS

\plancheta{Atividade de Som 40 --- Contagem de fonemas}

\subsection{Atividade de Som 40 --- Contagem de fonemas}
\begin{exercicio}[Atividade de Som 40 --- Contagem de fonemas]

\metadata{Som40}{Silábico-Alfabético}{EF02LP01, EF02LP04}{D1}{EF02LP04}{CNE/CEB nº 2/2012}

\kumontiming{2}{90}

\noindent\textbf{Nome:} \underline{\hspace{3.2cm}} \quad \textbf{Data:} \underline{\hspace{1.6cm}} \quad \textbf{Nota:} \underline{\hspace{0.9cm}}/10


\noindent\textbf{Comando:} Conte os fonemas (sons) de: CHAVE (4), PRATO (5), FLOR (4), CASA (4), QUEIJO (5). Compare com o número de letras.


\begin{center}
\begin{tikzpicture}
  \draw[thick, dashed] (0,0) rectangle (12,6);
\end{tikzpicture}
\end{center}

\noindent\textbf{Escreva 2 palavras com o som da atividade:}

\begin{center}
  \underline{\hspace{3cm}} \quad \underline{\hspace{3cm}}
\end{center}

\noindent\textbf{Desenhe uma situação com a palavra da atividade:}

\begin{center}
\begin{tikzpicture}
  \draw[thick, dashed] (0,0) rectangle (12,4.5);
\end{tikzpicture}
\end{center}

\end{exercicio}

%% ATIVIDADE 41: ORDEM DOS SONS

\plancheta{Atividade de Som 41 --- Ordem dos sons}

\subsection{Atividade de Som 41 --- Ordem dos sons}
\begin{exercicio}[Atividade de Som 41 --- Ordem dos sons]

\metadata{Som41}{Silábico-Alfabético}{EF02LP01, EF02LP04}{D1}{EF02LP04}{CNE/CEB nº 2/2012}

\kumontiming{2}{90}

\noindent\textbf{Nome:} \underline{\hspace{3.2cm}} \quad \textbf{Data:} \underline{\hspace{1.6cm}} \quad \textbf{Nota:} \underline{\hspace{0.9cm}}/10


\noindent\textbf{Comando:} Diga PRATO e responda: qual é o 1º som? (P) E o 2º? (R) E o 3º? (A). Faça o mesmo com TREM e FLOR.


\begin{center}
\begin{tikzpicture}
  \draw[thick, dashed] (0,0) rectangle (12,6);
\end{tikzpicture}
\end{center}

\noindent\textbf{Escreva 2 palavras com o som da atividade:}

\begin{center}
  \underline{\hspace{3cm}} \quad \underline{\hspace{3cm}}
\end{center}

\noindent\textbf{Desenhe uma situação com a palavra da atividade:}

\begin{center}
\begin{tikzpicture}
  \draw[thick, dashed] (0,0) rectangle (12,4.5);
\end{tikzpicture}
\end{center}

\end{exercicio}

%% ATIVIDADE 42: ASSOCIAÇÃO SOM-GRAFIA FINAL

\plancheta{Atividade de Som 42 --- Associação som-grafia final}

\subsection{Atividade de Som 42 --- Associação som-grafia final}
\begin{exercicio}[Atividade de Som 42 --- Associação som-grafia final]

\metadata{Som42}{Silábico-Alfabético}{EF02LP01, EF02LP04}{D1}{EF02LP04}{CNE/CEB nº 2/2012}

\kumontiming{2}{90}

\noindent\textbf{Nome:} \underline{\hspace{3.2cm}} \quad \textbf{Data:} \underline{\hspace{1.6cm}} \quad \textbf{Nota:} \underline{\hspace{0.9cm}}/10


\noindent\textbf{Comando:} Ouça o som e escolha a grafia: /ch/ pode ser CH ou X (chave, xícara). /s/ pode ser S, Ç, SS ou X (sapo, coração, passar, texto). Registre por escrito suas escolhas.


\begin{center}
\begin{tikzpicture}
  \draw[thick, dashed] (0,0) rectangle (12,6);
\end{tikzpicture}
\end{center}

\noindent\textbf{Escreva 2 palavras com o som da atividade:}

\begin{center}
  \underline{\hspace{3cm}} \quad \underline{\hspace{3cm}}
\end{center}

\noindent\textbf{Desenhe uma situação com a palavra da atividade:}

\begin{center}
\begin{tikzpicture}
  \draw[thick, dashed] (0,0) rectangle (12,4.5);
\end{tikzpicture}
\end{center}

\end{exercicio}

%% ATIVIDADE 51: SEGMENTAÇÃO DE FRASES EM PALAVRAS

\plancheta{Atividade de Som 51 --- Segmentação de frases em palavras}

\subsection{Atividade de Som 51 --- Segmentação de frases em palavras}
\begin{exercicio}[Atividade de Som 51 --- Segmentação de frases em palavras]

\metadata{Som51}{Silábico-Alfabético}{EF02LP01, EF02LP04}{D1}{EF02LP04}{CNE/CEB nº 2/2012}

\kumontiming{2}{90}

\noindent\textbf{Nome:} \underline{\hspace{3.2cm}} \quad \textbf{Data:} \underline{\hspace{1.6cm}} \quad \textbf{Nota:} \underline{\hspace{0.9cm}}/10


\noindent\textbf{Comando:} Ouça a frase 'O coelho fugiu do jardim' e bata palmas a cada palavra: 5 palavras. Faça com 3 frases do seu livro e escreva quantas palavras cada uma tem.


\begin{center}
\begin{tikzpicture}
  \draw[thick, dashed] (0,0) rectangle (12,6);
\end{tikzpicture}
\end{center}

\noindent\textbf{Escreva 2 palavras com o som da atividade:}

\begin{center}
  \underline{\hspace{3cm}} \quad \underline{\hspace{3cm}}
\end{center}

\noindent\textbf{Desenhe uma situação com a palavra da atividade:}

\begin{center}
\begin{tikzpicture}
  \draw[thick, dashed] (0,0) rectangle (12,4.5);
\end{tikzpicture}
\end{center}

\end{exercicio}

%% ATIVIDADE 52: DITADO DE FRASES DO ANO

\plancheta{Atividade de Som 52 --- Ditado de frases do ano}

\subsection{Atividade de Som 52 --- Ditado de frases do ano}
\begin{exercicio}[Atividade de Som 52 --- Ditado de frases do ano]

\metadata{Som52}{Silábico-Alfabético}{EF02LP01, EF02LP04}{D1}{EF02LP04}{CNE/CEB nº 2/2012}

\kumontiming{2}{90}

\noindent\textbf{Nome:} \underline{\hspace{3.2cm}} \quad \textbf{Data:} \underline{\hspace{1.6cm}} \quad \textbf{Nota:} \underline{\hspace{0.9cm}}/10


\noindent\textbf{Comando:} O professor dita 5 frases com as grafias do ano: 1) A chave abriu o portão. 2) O coelho comeu a cenoura. 3) O carro parou na chuva. 4) A casa fica perto da praça. 5) O coração bate forte. Escreva e corrija.


\begin{center}
\begin{tikzpicture}
  \draw[thick, dashed] (0,0) rectangle (12,6);
\end{tikzpicture}
\end{center}

\noindent\textbf{Escreva 2 palavras com o som da atividade:}

\begin{center}
  \underline{\hspace{3cm}} \quad \underline{\hspace{3cm}}
\end{center}

\noindent\textbf{Desenhe uma situação com a palavra da atividade:}

\begin{center}
\begin{tikzpicture}
  \draw[thick, dashed] (0,0) rectangle (12,4.5);
\end{tikzpicture}
\end{center}

\end{exercicio}

%% ATIVIDADE 53: REVISÃO GERAL — RELAXAMENTO ARTICULATÓRIO

\plancheta{Atividade de Som 53 --- Revisão Geral — Relaxamento Articulatório}

\subsection{Atividade de Som 53 --- Revisão Geral — Relaxamento Articulatório}
\begin{exercicio}[Atividade de Som 53 --- Revisão Geral — Relaxamento Articulatório]

\metadata{Som53}{Silábico-Alfabético}{EF02LP01, EF02LP04}{D1}{EF02LP04}{CNE/CEB nº 2/2012}

\kumontiming{2}{90}

\noindent\textbf{Nome:} \underline{\hspace{3.2cm}} \quad \textbf{Data:} \underline{\hspace{1.6cm}} \quad \textbf{Nota:} \underline{\hspace{0.9cm}}/10


\noindent\textbf{Comando:} O aluno revisa as boquinhas das grafias estudadas, compara as sílabas com CH/C e NH/LH, completa um quadro de revisão com as palavras-alvo e escreve uma frase usando duas grafias revisadas. Encerra com a leitura em voz alta de um texto curto com todas as grafias do ano.


\begin{center}
\begin{tikzpicture}
  \draw[thick, dashed] (0,0) rectangle (12,6);
\end{tikzpicture}
\end{center}

\noindent\textbf{Escreva 2 palavras com o som da atividade:}

\begin{center}
  \underline{\hspace{3cm}} \quad \underline{\hspace{3cm}}
\end{center}

\noindent\textbf{Desenhe uma situação com a palavra da atividade:}

\begin{center}
\begin{tikzpicture}
  \draw[thick, dashed] (0,0) rectangle (12,4.5);
\end{tikzpicture}
\end{center}

\end{exercicio}
